% Options for packages loaded elsewhere
\PassOptionsToPackage{unicode}{hyperref}
\PassOptionsToPackage{hyphens}{url}
\documentclass[
]{article}
\usepackage{xcolor}
\usepackage[margin=1in]{geometry}
\usepackage{amsmath,amssymb}
\setcounter{secnumdepth}{-\maxdimen} % remove section numbering
\usepackage{iftex}
\ifPDFTeX
  \usepackage[T1]{fontenc}
  \usepackage[utf8]{inputenc}
  \usepackage{textcomp} % provide euro and other symbols
\else % if luatex or xetex
  \usepackage{unicode-math} % this also loads fontspec
  \defaultfontfeatures{Scale=MatchLowercase}
  \defaultfontfeatures[\rmfamily]{Ligatures=TeX,Scale=1}
\fi
\usepackage{lmodern}
\ifPDFTeX\else
  % xetex/luatex font selection
\fi
% Use upquote if available, for straight quotes in verbatim environments
\IfFileExists{upquote.sty}{\usepackage{upquote}}{}
\IfFileExists{microtype.sty}{% use microtype if available
  \usepackage[]{microtype}
  \UseMicrotypeSet[protrusion]{basicmath} % disable protrusion for tt fonts
}{}
\makeatletter
\@ifundefined{KOMAClassName}{% if non-KOMA class
  \IfFileExists{parskip.sty}{%
    \usepackage{parskip}
  }{% else
    \setlength{\parindent}{0pt}
    \setlength{\parskip}{6pt plus 2pt minus 1pt}}
}{% if KOMA class
  \KOMAoptions{parskip=half}}
\makeatother
\usepackage{color}
\usepackage{fancyvrb}
\newcommand{\VerbBar}{|}
\newcommand{\VERB}{\Verb[commandchars=\\\{\}]}
\DefineVerbatimEnvironment{Highlighting}{Verbatim}{commandchars=\\\{\}}
% Add ',fontsize=\small' for more characters per line
\usepackage{framed}
\definecolor{shadecolor}{RGB}{248,248,248}
\newenvironment{Shaded}{\begin{snugshade}}{\end{snugshade}}
\newcommand{\AlertTok}[1]{\textcolor[rgb]{0.94,0.16,0.16}{#1}}
\newcommand{\AnnotationTok}[1]{\textcolor[rgb]{0.56,0.35,0.01}{\textbf{\textit{#1}}}}
\newcommand{\AttributeTok}[1]{\textcolor[rgb]{0.13,0.29,0.53}{#1}}
\newcommand{\BaseNTok}[1]{\textcolor[rgb]{0.00,0.00,0.81}{#1}}
\newcommand{\BuiltInTok}[1]{#1}
\newcommand{\CharTok}[1]{\textcolor[rgb]{0.31,0.60,0.02}{#1}}
\newcommand{\CommentTok}[1]{\textcolor[rgb]{0.56,0.35,0.01}{\textit{#1}}}
\newcommand{\CommentVarTok}[1]{\textcolor[rgb]{0.56,0.35,0.01}{\textbf{\textit{#1}}}}
\newcommand{\ConstantTok}[1]{\textcolor[rgb]{0.56,0.35,0.01}{#1}}
\newcommand{\ControlFlowTok}[1]{\textcolor[rgb]{0.13,0.29,0.53}{\textbf{#1}}}
\newcommand{\DataTypeTok}[1]{\textcolor[rgb]{0.13,0.29,0.53}{#1}}
\newcommand{\DecValTok}[1]{\textcolor[rgb]{0.00,0.00,0.81}{#1}}
\newcommand{\DocumentationTok}[1]{\textcolor[rgb]{0.56,0.35,0.01}{\textbf{\textit{#1}}}}
\newcommand{\ErrorTok}[1]{\textcolor[rgb]{0.64,0.00,0.00}{\textbf{#1}}}
\newcommand{\ExtensionTok}[1]{#1}
\newcommand{\FloatTok}[1]{\textcolor[rgb]{0.00,0.00,0.81}{#1}}
\newcommand{\FunctionTok}[1]{\textcolor[rgb]{0.13,0.29,0.53}{\textbf{#1}}}
\newcommand{\ImportTok}[1]{#1}
\newcommand{\InformationTok}[1]{\textcolor[rgb]{0.56,0.35,0.01}{\textbf{\textit{#1}}}}
\newcommand{\KeywordTok}[1]{\textcolor[rgb]{0.13,0.29,0.53}{\textbf{#1}}}
\newcommand{\NormalTok}[1]{#1}
\newcommand{\OperatorTok}[1]{\textcolor[rgb]{0.81,0.36,0.00}{\textbf{#1}}}
\newcommand{\OtherTok}[1]{\textcolor[rgb]{0.56,0.35,0.01}{#1}}
\newcommand{\PreprocessorTok}[1]{\textcolor[rgb]{0.56,0.35,0.01}{\textit{#1}}}
\newcommand{\RegionMarkerTok}[1]{#1}
\newcommand{\SpecialCharTok}[1]{\textcolor[rgb]{0.81,0.36,0.00}{\textbf{#1}}}
\newcommand{\SpecialStringTok}[1]{\textcolor[rgb]{0.31,0.60,0.02}{#1}}
\newcommand{\StringTok}[1]{\textcolor[rgb]{0.31,0.60,0.02}{#1}}
\newcommand{\VariableTok}[1]{\textcolor[rgb]{0.00,0.00,0.00}{#1}}
\newcommand{\VerbatimStringTok}[1]{\textcolor[rgb]{0.31,0.60,0.02}{#1}}
\newcommand{\WarningTok}[1]{\textcolor[rgb]{0.56,0.35,0.01}{\textbf{\textit{#1}}}}
\usepackage{graphicx}
\makeatletter
\newsavebox\pandoc@box
\newcommand*\pandocbounded[1]{% scales image to fit in text height/width
  \sbox\pandoc@box{#1}%
  \Gscale@div\@tempa{\textheight}{\dimexpr\ht\pandoc@box+\dp\pandoc@box\relax}%
  \Gscale@div\@tempb{\linewidth}{\wd\pandoc@box}%
  \ifdim\@tempb\p@<\@tempa\p@\let\@tempa\@tempb\fi% select the smaller of both
  \ifdim\@tempa\p@<\p@\scalebox{\@tempa}{\usebox\pandoc@box}%
  \else\usebox{\pandoc@box}%
  \fi%
}
% Set default figure placement to htbp
\def\fps@figure{htbp}
\makeatother
\setlength{\emergencystretch}{3em} % prevent overfull lines
\providecommand{\tightlist}{%
  \setlength{\itemsep}{0pt}\setlength{\parskip}{0pt}}
\usepackage{bookmark}
\IfFileExists{xurl.sty}{\usepackage{xurl}}{} % add URL line breaks if available
\urlstyle{same}
\hypersetup{
  pdftitle={Sample size{,} power{,} and Quarto reporting},
  hidelinks,
  pdfcreator={LaTeX via pandoc}}

\title{Sample size, power, and Quarto reporting}
\author{}
\date{\vspace{-2.5em}}

\begin{document}
\maketitle

\textbf{Workflow lab:} Goal → Approach → Execution → Check → Report.

\subsection{Learning objectives}\label{learning-objectives}

\begin{itemize}
\tightlist
\item
  Compute analytic sample size for the core inferential designs.
\item
  Produce the same answer by simulation and understand when each
  approach is preferable.
\item
  Assemble a short Quarto report whose numbers, tables, and figures all
  regenerate from a single render.
\end{itemize}

\subsection{Prerequisites}\label{prerequisites}

Courses 1 Weeks 1--4 up to this lab; \texttt{pwr}, \texttt{gtsummary},
and \texttt{broom} installed.

\subsection{Background}\label{background}

Sample size is one of the few places where applied statistics is
consulted \emph{before} the data exist. A study that is too small cannot
distinguish a real effect from noise; a study that is too large wastes
participants and money. The four ingredients of a sample-size
calculation --- effect size, variability, the significance level, and
power --- translate directly between a clinical protocol and a funding
application, and no competent ethics committee will approve a protocol
without them.

There are two complementary approaches. Closed-form formulae, packaged
in \texttt{pwr} and friends, are fast, transparent, and correct for the
textbook designs: one- and two-sample t-tests, one- and two-proportion
tests, correlation, and ANOVA. Simulation-based power, by contrast, is
the honest answer for anything the textbooks skip --- mixed models,
skewed outcomes, interim analyses --- and costs only a few minutes of
compute. The habit to build is to compute analytically first, then
verify by simulation; if the two disagree, either the formula does not
apply to your design or your simulation is wrong, and you need to know
which.

Reporting is the other half of this lab. Every number in a well-written
report should trace back to the code that produced it. Quarto, with
\texttt{gtsummary} for tables and \texttt{broom::tidy()} for model
output, makes this nearly automatic.

\subsection{Setup}\label{setup}

\begin{Shaded}
\begin{Highlighting}[]
\FunctionTok{library}\NormalTok{(tidyverse)}
\FunctionTok{library}\NormalTok{(pwr)}
\FunctionTok{library}\NormalTok{(gtsummary)}
\FunctionTok{library}\NormalTok{(broom)}
\FunctionTok{set.seed}\NormalTok{(}\DecValTok{42}\NormalTok{)}
\FunctionTok{theme\_set}\NormalTok{(}\FunctionTok{theme\_minimal}\NormalTok{(}\AttributeTok{base\_size =} \DecValTok{12}\NormalTok{))}
\end{Highlighting}
\end{Shaded}

\subsection{1. Goal}\label{goal}

Plan a hypothetical two-arm trial comparing a new antihypertensive with
placebo, choose a sample size with 80\% power to detect a 5 mmHg
difference at α = 0.05, verify by simulation, and produce a reporting
table we can paste into the protocol.

\subsection{2. Approach}\label{approach}

Closed-form first. Cohen's \emph{d} for a 5 mmHg effect on a scale with
SD 10 mmHg is \emph{d} = 0.5 --- a medium effect by convention.

\begin{Shaded}
\begin{Highlighting}[]
\NormalTok{  d }\OtherTok{=} \FloatTok{0.5}\NormalTok{, power }\OtherTok{=} \FloatTok{0.80}\NormalTok{, sig.level }\OtherTok{=} \FloatTok{0.05}\NormalTok{,}
\NormalTok{  type }\OtherTok{=} \StringTok{"two.sample"}\NormalTok{, alternative }\OtherTok{=} \StringTok{"two.sided"}
\ErrorTok{)}
\NormalTok{pwr\_calc}
\end{Highlighting}
\end{Shaded}

The formula calls for about \texttt{ceiling(pwr\_calc\$n)} per arm.

\begin{Shaded}
\begin{Highlighting}[]
\NormalTok{ns }\OtherTok{\textless{}{-}} \FunctionTok{seq}\NormalTok{(}\DecValTok{10}\NormalTok{, }\DecValTok{200}\NormalTok{, }\AttributeTok{by =} \DecValTok{5}\NormalTok{)}
\NormalTok{powers }\OtherTok{\textless{}{-}} \FunctionTok{sapply}\NormalTok{(ns, }\ControlFlowTok{function}\NormalTok{(n) \{}
  \FunctionTok{pwr.t.test}\NormalTok{(}\AttributeTok{n =}\NormalTok{ n, }\AttributeTok{d =} \FloatTok{0.5}\NormalTok{, }\AttributeTok{sig.level =} \FloatTok{0.05}\NormalTok{,}
             \AttributeTok{type =} \StringTok{"two.sample"}\NormalTok{)}\SpecialCharTok{$}\NormalTok{power}
\NormalTok{\})}
\FunctionTok{tibble}\NormalTok{(}\AttributeTok{n\_per\_arm =}\NormalTok{ ns, }\AttributeTok{power =}\NormalTok{ powers) }\SpecialCharTok{|\textgreater{}}
  \FunctionTok{ggplot}\NormalTok{(}\FunctionTok{aes}\NormalTok{(n\_per\_arm, power)) }\SpecialCharTok{+}
  \FunctionTok{geom\_line}\NormalTok{(}\AttributeTok{linewidth =} \FloatTok{0.8}\NormalTok{) }\SpecialCharTok{+}
  \FunctionTok{geom\_hline}\NormalTok{(}\AttributeTok{yintercept =} \FloatTok{0.8}\NormalTok{, }\AttributeTok{linetype =} \DecValTok{2}\NormalTok{, }\AttributeTok{colour =} \StringTok{"grey50"}\NormalTok{) }\SpecialCharTok{+}
  \FunctionTok{labs}\NormalTok{(}\AttributeTok{x =} \StringTok{"N per arm"}\NormalTok{, }\AttributeTok{y =} \StringTok{"Power"}\NormalTok{,}
       \AttributeTok{title =} \StringTok{"Power curve for d = 0.5, α = 0.05"}\NormalTok{)}
\end{Highlighting}
\end{Shaded}

\subsection{3. Execution}\label{execution}

Verify by simulation. A function that runs the trial once and returns
the p-value:

\begin{Shaded}
\begin{Highlighting}[]
\NormalTok{  placebo }\OtherTok{\textless{}{-}} \FunctionTok{rnorm}\NormalTok{(n\_per\_arm, }\DecValTok{0}\NormalTok{, sd)}
\NormalTok{  drug    }\OtherTok{\textless{}{-}} \FunctionTok{rnorm}\NormalTok{(n\_per\_arm, }\SpecialCharTok{{-}}\NormalTok{delta, sd)}
  \FunctionTok{t.test}\NormalTok{(drug, placebo, }\AttributeTok{var.equal =} \ConstantTok{FALSE}\NormalTok{)}\SpecialCharTok{$}\NormalTok{p.value}
\ErrorTok{\}}

\NormalTok{sim\_power }\OtherTok{\textless{}{-}} \ControlFlowTok{function}\NormalTok{(n\_per\_arm, }\AttributeTok{reps =} \DecValTok{500}\NormalTok{, ...) \{}
  \FunctionTok{mean}\NormalTok{(}\FunctionTok{replicate}\NormalTok{(reps, }\FunctionTok{sim\_trial}\NormalTok{(n\_per\_arm, ...)) }\SpecialCharTok{\textless{}} \FloatTok{0.05}\NormalTok{)}
\NormalTok{\}}

\NormalTok{n\_try  }\OtherTok{\textless{}{-}} \FunctionTok{ceiling}\NormalTok{(pwr\_calc}\SpecialCharTok{$}\NormalTok{n)}
\FunctionTok{sim\_power}\NormalTok{(n\_try, }\AttributeTok{reps =} \DecValTok{1000}\NormalTok{)}
\end{Highlighting}
\end{Shaded}

Within Monte-Carlo error of the analytic 0.80, as expected.

\subsection{4. Check}\label{check}

Now imagine the data have been collected. We simulate one realisation
and fit the planned analysis.

\begin{Shaded}
\begin{Highlighting}[]
\NormalTok{trial }\OtherTok{\textless{}{-}} \FunctionTok{tibble}\NormalTok{(}
  \AttributeTok{arm =} \FunctionTok{rep}\NormalTok{(}\FunctionTok{c}\NormalTok{(}\StringTok{"placebo"}\NormalTok{, }\StringTok{"drug"}\NormalTok{), }\AttributeTok{each =}\NormalTok{ n),}
  \AttributeTok{sbp\_change =} \FunctionTok{c}\NormalTok{(}\FunctionTok{rnorm}\NormalTok{(n, }\DecValTok{0}\NormalTok{, }\DecValTok{10}\NormalTok{), }\FunctionTok{rnorm}\NormalTok{(n, }\SpecialCharTok{{-}}\DecValTok{5}\NormalTok{, }\DecValTok{10}\NormalTok{))}
\NormalTok{)}

\NormalTok{fit }\OtherTok{\textless{}{-}} \FunctionTok{t.test}\NormalTok{(sbp\_change }\SpecialCharTok{\textasciitilde{}}\NormalTok{ arm, }\AttributeTok{data =}\NormalTok{ trial, }\AttributeTok{var.equal =} \ConstantTok{FALSE}\NormalTok{)}
\NormalTok{broom}\SpecialCharTok{::}\FunctionTok{tidy}\NormalTok{(fit)}
\end{Highlighting}
\end{Shaded}

\subsection{5. Report}\label{report}

\begin{Shaded}
\begin{Highlighting}[]
\NormalTok{trial }\SpecialCharTok{|\textgreater{}}
  \FunctionTok{tbl\_summary}\NormalTok{(}
    \AttributeTok{by =}\NormalTok{ arm,}
    \AttributeTok{statistic =} \FunctionTok{list}\NormalTok{(}\FunctionTok{all\_continuous}\NormalTok{() }\SpecialCharTok{\textasciitilde{}} \StringTok{"\{mean\} (\{sd\})"}\NormalTok{),}
    \AttributeTok{digits    =} \FunctionTok{all\_continuous}\NormalTok{() }\SpecialCharTok{\textasciitilde{}} \DecValTok{1}\NormalTok{,}
    \AttributeTok{label     =} \FunctionTok{list}\NormalTok{(}\AttributeTok{sbp\_change =} \StringTok{"Change in SBP (mmHg)"}\NormalTok{)}
\NormalTok{  ) }\SpecialCharTok{|\textgreater{}}
  \FunctionTok{add\_p}\NormalTok{() }\SpecialCharTok{|\textgreater{}}
  \FunctionTok{modify\_caption}\NormalTok{(}\StringTok{"**Table 1. Change in systolic blood pressure, by arm.**"}\NormalTok{)}
\end{Highlighting}
\end{Shaded}

\begin{quote}
In a simulated two-arm trial (\emph{n} = \texttt{n} per arm), mean
change in systolic blood pressure was lower in the drug arm than in
placebo by \texttt{round(-diff(fit\$estimate),\ 1)} mmHg (95\% CI:
\texttt{round(-fit\$conf.int{[}2{]},\ 1)} to
\texttt{round(-fit\$conf.int{[}1{]},\ 1)}; \emph{p} =
\texttt{signif(fit\$p.value,\ 2)}).
\end{quote}

\subsection{Common pitfalls}\label{common-pitfalls}

\begin{itemize}
\tightlist
\item
  Plugging a standardised effect size into a sample-size calculator
  without stopping to think whether it is plausible in your field.
\item
  Quoting ``power = 0.8'' as a universal target when your decision would
  justify 0.9 or 0.95.
\item
  Reporting a \emph{p}-value without the effect size and interval that
  make it interpretable.
\item
  Pasting numbers into the manuscript by hand --- anything that breaks
  the render trail will eventually break the numbers.
\end{itemize}

\subsection{Further reading}\label{further-reading}

\begin{itemize}
\tightlist
\item
  \href{../../appendices/D-writing-a-report.Rmd}{Writing a report}.
\item
  Champely S. (2024). \emph{pwr: Basic functions for power analysis.}
\item
  Cohen J. (1988). \emph{Statistical Power Analysis for the Behavioral
  Sciences.}
\end{itemize}

\subsection{Session info}\label{session-info}

\begin{Shaded}
\begin{Highlighting}[]

\end{Highlighting}
\end{Shaded}

\subsection{See also --- chapter index}\label{see-also-chapter-index}

\begin{itemize}
\tightlist
\item
  \href{../index.Rmd}{Methods in this chapter}
\item
  \href{../../../ch00-find-your-method.Rmd}{Find your method (Chapter
  0)}
\end{itemize}

\end{document}
